Finally, as a result, interaction between designs now has unambiguous meaning. There are meaningful implications for securing systems because the intention of the interaction is unambiguous. While before, the mechanical sequence had to be correct, now the design can ensure that it is the right narrative that is communicating with it. If the design is looked at as the brain of the software system, then this ensures that the correct brain is speaking, making it more difficult to compromise a system.

There are many approaches to implement this but one could be that the DAL can be extended to generate a dynamic checksum along the narrative path to identify unique narrative identity. There are more interesting ways that I can think of for securing it further by sharing the identity of the instance of the brain and using that to generate the proof. Regardless, the point here is, by eliminating ambiguity in a meaningful way, more automation will be enabled and securing the interaction between and within designs is another application for this framework.